% ID: SEL-D0067-OA-E03
% ESTADO: cotejada. Segunda pasada visual y borrador compilado revisados.
% Fuente: Selectividad/sel_2009_matematicas_aplicadas/mat apli CC SS01.pdf
% SHA-256: f5e34036805266809f1b5a2c6087d7ae598afd929b5179ec5ea3888db6cbccf4
% Página: 1; opción A; ejercicio original 3.

\noindent\underline{Parte I}

Lena y Adrián son aficionados al tiro con arco. Lena da en el blanco con probabilidad $\dfrac{7}{11}$ y Adrián con probabilidad $\dfrac{9}{13}$. Si ambos sucesos son independientes, calcule la probabilidad de los siguientes sucesos:
\begin{enumerate}[label=\alph*)]
\item \textbf{(0.6 puntos)} ``Ambos dan en el blanco''.
\item \textbf{(0.6 puntos)} ``Sólo Lena da en el blanco''.
\item \textbf{(0.8 puntos)} ``Al menos uno da en el blanco''.
\end{enumerate}
\begin{samepage}
\noindent\underline{Parte II}

En una muestra aleatoria de 100 individuos se ha obtenido, para la edad, una media de $17.5$ años. Se sabe que la edad en la población, de la que procede esa muestra, sigue una distribución Normal con una desviación típica de $0.8$ años.
\begin{enumerate}[label=\alph*)]
\item \textbf{(1.5 puntos)} Obtenga un intervalo de confianza, al $94\%$, para la edad media de la población.
\item \textbf{(0.5 puntos)} \textquestiondown{}Qué error máximo se comete en la estimación anterior?
\end{enumerate}
\end{samepage}
